\documentclass[a4paper,10pt,landscape]{article}

\usepackage[czech]{babel}
\usepackage[T1]{fontenc}
\usepackage[utf8]{inputenc}
\usepackage{amsmath,amssymb}
\usepackage{geometry}
\usepackage{enumitem}

\geometry{a4paper,landscape,left=8mm,right=8mm,top=8mm,bottom=8mm}
\pagestyle{empty}
\setlength{\parindent}{0pt}
\setlength{\parskip}{0pt}
\setlist[enumerate]{leftmargin=*,itemsep=2pt,topsep=2pt}

% Uprav název testu; jednotlivé varianty jsou vždy A/B, C/D atd.
\newcommand{\TestTitle}{Komplexní čísla}
\newcommand{\TestHeader}[1]{%
    {\bfseries TEST -- \TestTitle}\hfill{\bfseries Varianta #1}\\[2mm]
    Jméno a příjmení: \rule{4.5cm}{0.4pt}\hfill Třída: \rule{1.5cm}{0.4pt}\\[2mm]
    Datum: \rule{2.4cm}{0.4pt}\hfill Počet bodů: \rule{1.4cm}{0.4pt}\hfill Známka: \rule{1.4cm}{0.4pt}\\[1mm]
    \hrule\vspace{3mm}
}
\newcommand{\AnswerSpace}[1]{\vspace{#1}}

\begin{document}
\noindent
\begin{minipage}[t]{0.48\textwidth}
    \small
    \TestHeader{A}

    % ZADÁNÍ VARIANTY A
    \begin{enumerate}
        \item Upravte do algebraického tvaru: \[3i^{14}-2i^{19}+i^{24}-4i^7\]
        \AnswerSpace{10mm}
        \item Převeďte do goniometrického tvaru: \[z=-\sqrt3+i\]
        \AnswerSpace{12mm}
        \item Převeďte do algebraického tvaru: \[z=6(\cos150^\circ+i\sin150^\circ)\]
        \AnswerSpace{10mm}
        \item Jsou dána komplexní čísla \(z_1=3+2i\) a \(z_2=1-i\). Vypočítejte: \[z_1+z_2,\quad z_1-z_2,\quad z_1z_2,\quad \frac{z_1}{z_2}\]
        \AnswerSpace{25mm}
        \item Jsou dána komplexní čísla: \[\begin{aligned}z_1&=4(\cos120^\circ+i\sin120^\circ)\\z_2&=2(\cos30^\circ+i\sin30^\circ)\end{aligned}\] Určete v goniometrickém tvaru součin \(z_1z_2\) a podíl \(\frac{z_1}{z_2}\).
        \AnswerSpace{20mm}
    \end{enumerate}
\end{minipage}%
\hfill\vrule width 0.4pt\hfill
\begin{minipage}[t]{0.48\textwidth}
    \small
    \TestHeader{B}

    % ZADÁNÍ VARIANTY B
    \begin{enumerate}
        \item Upravte do algebraického tvaru: \[2i^{15}+3i^{18}-i^{21}+5i^{28}\]
        \AnswerSpace{10mm}
        \item Převeďte do goniometrického tvaru: \[z=-1+\sqrt3i\]
        \AnswerSpace{12mm}
        \item Převeďte do algebraického tvaru: \[z=8(\cos210^\circ+i\sin210^\circ)\]
        \AnswerSpace{10mm}
        \item Jsou dána komplexní čísla \(z_1=2+3i\) a \(z_2=1+i\). Vypočítejte: \[z_1+z_2,\quad z_1-z_2,\quad z_1z_2,\quad \frac{z_1}{z_2}\]
        \AnswerSpace{25mm}
        \item Jsou dána komplexní čísla: \[\begin{aligned}z_1&=6(\cos150^\circ+i\sin150^\circ)\\z_2&=3(\cos60^\circ+i\sin60^\circ)\end{aligned}\] Určete v goniometrickém tvaru součin \(z_1z_2\) a podíl \(\frac{z_1}{z_2}\).
        \AnswerSpace{20mm}
    \end{enumerate}
\end{minipage}
\end{document}
